\documentclass{article}
\usepackage{graphicx}
\usepackage[margin=1in, top=0.65in]{geometry}
\usepackage{enumitem}
\usepackage{titlesec}
\usepackage{parskip}
\usepackage{hyperref}
\usepackage{amsmath,amssymb}

\hypersetup{colorlinks=true, urlcolor=blue, linkcolor=blue}

\titleformat{\section}{\large\bfseries}{}{0em}{}[\titlerule]
\titleformat{\subsection}{\normalsize\bfseries}{}{0em}{}
\titlespacing*{\section}{0pt}{6pt}{4pt}

\setlist[enumerate]{noitemsep, topsep=2pt}
\setlist[itemize]{noitemsep, topsep=2pt}

\title{ASEN 5264 Proposal}
\author{Allen Devaraj Augustin Ponraj}
\date{March 2026}

\begin{document}
\maketitle
\vspace{-12pt}

\section{Title \& Goal}

\textbf{Robot Pick-and-Place Under Observation Uncertainty.} I propose to compare three decision-making approaches for a simulated robot arm performing pick-and-place of a target block under \emph{observation uncertainty}: (1)~a baseline policy that sees only noisy observations; (2)~a belief-augmented policy that uses a particle filter and learns when to gather information vs.\ act; and (3)~a stretch goal using POMCP with a learned world model. The setting is a POMDP: the true block pose is hidden; the agent receives noisy or intermittent pose estimates and pays a cost per timestep (cost of looking), creating a clear exploration--exploitation trade-off.

\section{Uncertainty Model}

I inject three forms of observation uncertainty (no camera hardware model by default):

\begin{itemize}
  \item \textbf{Episodic noise:} Each episode gets a different observation noise level $\sigma_{\text{ep}}$ so some episodes are easy and some hard.
  \item \textbf{Multi-block occlusion:} A distractor block can occlude the target (as lego blocks are spawned in random places for each episode) from the camera so the agent sometimes receives no pose update (or stale).
  \item \textbf{Cost of looking:} A step penalty so that gathering information (e.g., moving to un-occlude or waiting for more observations) trades off with acting.
\end{itemize}

\section{Methods}

\begin{enumerate}
  \item \textbf{Plain PPO (baseline):} The agent observes raw noisy pose estimates and joint state. I train with PPO (e.g., Stable-Baselines3) to establish a ceiling for ``ignoring uncertainty.''
  \item \textbf{Belief-Augmented PPO (main):} A particle filter maintains a belief over block pose; the agent observes belief summary statistics (mean and variance). Same PPO architecture and reward. The agent learns \emph{when} to look (reduce uncertainty) vs.\ when to grasp (exploit), i.e., exploration-exploitation under cost of looking.
  \item \textbf{POMCP with learned world model (stretch):} Use a learned transition model (trained on data from the belief-augmented policy) to run POMCP tree search over the belief state, with the same action space and reward.
\end{enumerate}

\section{Task \& Metrics}

Task: pick the target block and place it at a random goal; max steps per episode with step penalty. Success is placement within a threshold of the goal. Metrics: success rate; mean steps to success (efficiency); steps and $\sigma$ at grasp time (whether the agent waited for confidence); success vs.\ $\sigma_{\text{ep}}$ sweep (robustness).

\section{Deliverables}

A Gymnasium-compatible pick-and-place environment with the above uncertainty sources; trained policies for the baseline and belief-augmented methods (and optionally POMCP); comparison plots (learning curves, success vs.\ noise, $\sigma$ at grasp); short video demos; and a final report summarizing the comparison and lessons for decision making under uncertainty.

\end{document}
